\chapter*{\Huge Conclusions}
\thispagestyle{thesis} % Force the thesis page style
\markboth{\MakeMarkcase{Conclusions}}{\MakeMarkcase{Conclusions}} % Set headers
\addcontentsline{toc}{chapter}{Conclusions} % Add to Table of Contents

\begin{enumerate}
\item Water-saving strategies, are efficient climate change mitigation practices, when compared to continuous flooding, but can result in negative effects on biodiversity conservation and rice production goals. Their mitigation potential is positively correlated to the strategy's intensity along the water use gradient (i.e., AWD being the highest intensity strategy and having the highest mitigation potential, followed by MSD). This correlation is inverted for the abundance of aquatic organisms in paddy fields, with AWD resulting in the lowest abundances and MSD buffering this effect. Yields are decreased significantly under AWD, while MSD allows similar production levels as in continuously flooded fields.\\ 

\item There is an important time component in water-saving strategies effect on CH$_{4}$ emissions, as fallow season emissions are affected by their their implementation during the growing season. Mitigation effects observed during growing seasons are inverted in fallow seasons, being AWD the strategy with highest fallow and overall annual emissions. Considering both seasons, MSD achieves the lowest cumulative CH$_{4}$ emissions. Such effects can be attributed to water-saving legacy effects on the structure and composition of soil microbial communities, hence modifying dominant metabolic pathways in rice paddy soils.\\

\item None of the assessed irrigation strategies is overall positive from an ecosystem multifunctionality perspective, as their effects are heterogeneous among indicators and agroecosystem services. Although overall multifuntionality is not affected, MSD and AWD have conflicting effects on the abundance and diversity of several organisms, and climate change mitigation and pest control indicators. Nevertheless, their strong positive effect decreasing CH$_{4}$ emissions drives a general positive effect from a climate change mitigation perspective. There are no evident links among different levels of agroecosystem services. Provisioning service, in the form of rice grain yields, is the only agroecosystem service affected significantly by water-saving strategies, being decreased under AWD.\\

% \item The implementation of water-saving irrigation practices must consider spatiotemporal dimensions specific to the assessed systems. From a spatial perspective, management goals are defined by landscape complexity... 

\item Assessing ecosystem effects of implementing water-saving strategies in rice paddy fields must follow a holistic perspective. Evaluation frameworks should account for potential ecological trade-offs between climate change mitigation and biodiversity conservation, effects on temporal CH$_{4}$ emission patterns for both growing and fallow seasons, and agroecosystem multifunctionality. Narrow assessments risk the promotion of strategies that might result detrimental on large temporal and spatial scales (e.g., implementing AWD only considering its climate change mitigation potential). Such holistic rationale must transcend rice paddies, applying across systems and management practices, providing decision-makers with the right tools to tackle global change and sustainable development challenges.\\

\end{enumerate}